% =========================================================
% Integrals - Techniques
% =========================================================
\subsection*{Techniques}

\begin{center}
\begin{tabular}{ll}
\toprule
Technique & Use when \\
\midrule
Substitution & The integrand contains a composition and its derivative \\
Integration by parts & The integrand is a product with a simpler derivative \\
Trigonometric integrals & Powers of trig functions can be reduced \\
Trigonometric substitution & Radicals match \(a^2-x^2\), \(a^2+x^2\), or \(x^2-a^2\) \\
Partial fractions & A rational function has a factorable denominator \\
Improper integrals & The interval is infinite or the integrand is unbounded \\
\bottomrule
\end{tabular}
\end{center}

\begin{remark*}[Method]
Before choosing a technique, simplify algebraically, check for a direct
antiderivative, and look for an inside function whose derivative is also
present.
\end{remark*}
